\section{Conclusiones}
Los cálculos permitieron reconstruir poblaciones de equilibrio condicionadas a masas de reposo, inventario HRG, temperatura y exceso cosmológico. A la referencia bibliográfica de 50 MeV se obtuvieron $2.7261\times10^{30}$ bariones y $2.7172\times10^{30}$ antibariones por cm$^3$, con exceso de $8.8765\times10^{27}$ cm$^{-3}$. Coexistieron 306.1147 antibariones por unidad del exceso neto; los nucleones y sus antipartículas aportaron 95.2176\% del total.

El criterio $n_{\bar B}=n_B-n_{\bar B}$ se alcanzó a 38.203739 MeV con PDG2025 y redondeó al hito publicado de 38.2 MeV. La variación de $\eta$ por sí sola desplazó el cruce a 38.192814--38.214598 MeV. El contraste aporta consistencia de modelos de equilibrio con condiciones relacionadas, sin determinar una incertidumbre física total.

La abundancia actual sola no identificó una población inicial única ni un historial de aniquilación. La aproximación $\mu_Q=\mu_S=0$ dejó carga y extrañeza netas distintas de cero. Estos límites delimitan la conclusión: se logró una reconstrucción térmica cuantitativa y reproducible, y una reconstrucción cosmológica completa requiere balances químicos y evolución cinética adicionales.

\section*{Disponibilidad de datos y código}
El paquete adjunto incluye fuentes de los manuscritos, figuras por idioma, catálogos fijados y licencia, 18 estados agregados para tres listas, 1329 registros por especie de partícula para PDG2025, curva del cruce, sensibilidad a $\eta$, 24 controles numéricos y guía de ejecución en Windows. Las conversiones cosmológicas y el análisis previo se conservan como antecedentes identificados. El cálculo se ejecutó en Python con una implementación independiente; se utilizaron los datos de Thermal-FIST, sin ejecutar su programa C++.

\section*{Asistencia en la preparación}
Se utilizó asistencia de inteligencia artificial para preparar redacción, código, figuras y maquetación. Los cálculos y controles descritos se ejecutaron y sus salidas se conservaron. El lanzador PowerShell se revisó estáticamente; no se ejecutó en Windows. La interpretación científica y la decisión de publicación corresponden al autor.
